% EM structure and teaching revision 20260928
\subsection{接地、孤立与相连导体的条件选择}
导体内部场为零只给出导体等势，还不足以决定它的电势值或总电荷。孤立导体给定总电荷 $Q$，其电势待求；接地导体给定电势，电荷可以经地线流入流出；导线连接的多个导体电势相等，但通常只守恒它们的总电荷。
例如两个相距很远的金属球，半径为 $R_1,R_2$，忽略相互影响。用细导线连接后
\[
 \frac{Q_1}{R_1}=\frac{Q_2}{R_2},\qquad Q_1+Q_2=Q,
 \qquad Q_i=\frac{QR_i}{R_1+R_2}.
\]
球面场强却满足 $E_1/E_2=R_2/R_1$；等势不意味着等电荷或等场强。若两球靠得很近，各自电势还包含对方电荷的贡献，上述孤立球近似便不再成立。
在有空腔的问题中，先用导体材料内部的高斯面确定内表面总感应电荷，再按接地或孤立条件确定外表面电荷。偏心电荷通常使内表面分布不均匀，却不妨碍用高斯定律求其总量。把“总电荷已知”与“分布已知”分开，是避免误用球对称公式的关键。

